Several studies have explored machine learning approaches for estimating parameters of sound synthesizers. The central choice in this work is to use Convolutional Neural Networks (CNNs) directly on raw audio signals. This choice is motivated by three practical aspects: raw waveforms preserve temporal structure, CNNs are well suited to high-dimensional inputs, and once trained they can provide fast inference. In addition, as emphasized by the ``Bitter Lesson'', greater reliance on data and computation often outperforms approaches based on hand-crafted representations~\cite{sutton2019bitterlesson}. In this sense, the present work follows a data-driven perspective in which the model learns the relevant patterns directly from the signal, without relying on a manual feature pipeline.

Claesson et al.~\cite{claesson2021resynthesis} proposed the re-synthesis of instrumental sounds using machine learning models combined with an FM synthesizer. Their approach, however, used STFT spectrograms as input rather than raw audio. Although this frequency-domain representation may simplify learning, it also introduces an additional preprocessing stage and does not fully exploit the temporal information present in the waveform. Steinmetz et al.~\cite{steinmetz2022ddx7} introduced DDX7, a differentiable FM synthesizer designed for musical instrument re-synthesis. Their method employs a Temporal Convolutional Network (TCN) and integrates Mel-spectrum losses, which provides a perceptually aligned training objective. The main trade-off is that the synthesizer must be embedded in the training loop, increasing computational and memory cost.

Finally, Kiranyaz et al.~\cite{kiranyaz2021survey} presented a comprehensive survey on one-dimensional CNNs and their applications, highlighting their effectiveness for sequential data such as audio signals. This survey supports the suitability of CNN-based architectures for learning temporal patterns directly from raw audio. In contrast to the works above, the present study emphasizes generalization through large-scale randomized exploration of the FM parameter space, while keeping the synthesizer outside the optimization process. This also makes the resulting model suitable for fast inference after training, which is important for interactive and real-time-oriented audio applications.
